% Generated draft: requires mathematical and source review.
\documentclass[11pt]{article}
\usepackage[margin=25mm]{geometry}
\usepackage{fontspec}
\setmainfont{Latin Modern Roman}
\usepackage{amsmath,amssymb,mathtools,unicode-math}
\setmathfont{Latin Modern Math}
\usepackage{enumitem,needspace,xurl,hyperref}
\setlength{\parindent}{0pt}
\setlength{\parskip}{5pt}
\setlength{\emergencystretch}{3em}
\newcommand{\N}{\mathbb{N}}
\newcommand{\Z}{\mathbb{Z}}
\newcommand{\Q}{\mathbb{Q}}
\newcommand{\R}{\mathbb{R}}
\newcommand{\C}{\mathbb{C}}
\newcommand{\F}{\mathbb{F}}
\newcommand{\A}{\mathbb{A}}
\newcommand{\PP}{\mathbb{P}}
\newcommand{\E}{\mathbb{E}}
\newcommand{\eps}{\varepsilon}
\newcommand{\dd}{\,\mathrm{d}}
\begin{document}
\section*{UnsolvedMath ID: 5500051}
\textbf{Source problem code:} catalogue entry
\subsection{Definitions and mathematical statement}
For a finite simple graph $G=(V,E)$, a three-dimensional straight-line grid drawing is an injection $p:V\to\Z^3$ such that each edge $uv$ is represented by the segment $[p(u),p(v)]$, no edge passes through a nonincident vertex, and two distinct edges intersect only at the image of a common endpoint, if they have one. The simple-graph convention is made explicit because loops or parallel edges need separate drawing rules.

For $n=|V|\ge1$ write, for $i=1,2,3$,
\[
 L_i=1+\max_{v\in V}p_i(v)-\min_{v\in V}p_i(v).
\]
A standard grid-box volume convention is $\operatorname{vol}(p)=L_1L_2L_3$, counting the integer-grid positions in the smallest axis-aligned box. The source says volume of the bounding box without clarifying this convention versus the product of geometric side lengths; the distinction matters for a flat drawing.

The question, under the grid-box convention, is whether
\[
 \exists C>0\ \forall n\ge1\ \forall G\text{ planar with }|V(G)|=n,
 \quad\exists p\text{ a 3D grid drawing of }G,
 \qquad\operatorname{vol}(p)\le Cn.
\]
The constant must be independent of $G$ and $n$. No prescribed aspect ratio or coordinate orientation beyond the integer grid is required.
\subsection{Short English statement}
Can every planar graph be drawn with straight, noncrossing edges at integer points in three dimensions inside a box of linear volume?

The registry and research summary label the question solved, asserting a positive answer via bounded queue number and related graph-layout results. The summary names Dujmović, Morin, and Wood, cites \emph{Layout of Graphs with Bounded Tree-Width}, and gives a 2015--2016 chronology with later constant improvements. Those attributions and dates are not verified here, and a bounded-tree-width result alone does not specify the required all-planar theorem. The background independently records the maintained TOPP entry as Open on access in 2026. This status conflict and the exact supporting theorem require review; neither label is silently selected as verified fact.
\subsection{Sources}
\begin{itemize}
\item Local UnsolvedMath registry, permanent ID 5500051, source code catalogue entry, ``Linear-Volume 3D Grid Drawings of Planar Graphs.'' Dataset curation and metadata: Ulam AI and UnsolvedMath Contributors, CC BY 4.0.
\item \emph{The Open Problems Project}, Computational Geometry, Problem 51, \url{https://topp.openproblem.net/p51}; source-TeX extraction, accessed 2026-07-29. Background records the maintained entry as Open.
\item Dujmović, Morin, and Wood, \emph{Layout of Graphs with Bounded Tree-Width}, as named in the supplied research summary; complete bibliographic information is absent. The summary also refers to Dujmović (2015), a 2015/2016 sequence, and subsequent constant improvements without exact papers or theorem statements.
\item Public source specification: \url{https://github.com/edemaine/topp}. The supplied rights note says its repository licence expressly covers software; problem-text reuse remains NEEDS\_REVIEW, separately from the dataset licence.
\end{itemize}
\end{document}
